\section{Nothing}

Why is there a base at all, why anything rather than nothing? The model gives a clean answer. There cannot be nothing.

Take the word strictly. True nothing is not empty space, which already has size, directions, and a clock. It is the absence of every distinction, everything the same as everything with no everything to speak of. Three things break it. It is a balance and not a resting place, a ball poised on a hilltop that the least difference sends rolling, with nothing to push it back. It contradicts itself the moment it is named, since to call it one whole, set apart from what it is not, is already to draw the very line it was meant to lack. And a whole that is genuinely everything has nothing else to relate to, so it can only relate to itself, and self-relation splits it in two, a side that notes and a side that is noted. That split is the first distinction.

From that one forced act the base is born. The two sides are the first two vibes, the relation between them the first connection, the side each takes the first tone, the act itself the first beat. Nothing came from outside, because there is no outside. The world bootstraps, lifting itself by its own self-relation, because holding still was never an option it had.

And there is no way back. To erase a distinction is itself a distinction, the difference between before and after, so distinctions only ever accumulate, and that one-way pile of what has happened is time. Something is not the strange case that needs explaining. Nothing is the one state the world cannot be in, and beginning was the only thing it could do.
